%% kapa: exam=1 round=37 year=2026 period=1 form=A subject=economics subject_ko=경제학원론 seq=39 number=79
%% source: Q-Net 감정평가사 기출문제 (cst00309) artlSeq=5258924 file=raw/1008/5258924/2026년 제37회 감정평가사 1차 1교시.pdf
%% conversion: kordoc 4.24.51 -> tools/gen1.py (auto); review=proofread-sonnet; flags=none
\begin{question}{79}
자본이동이 완전한 소국개방경제의 먼델-플레밍 모형에서 정부지출을 증가시킬 때 나타나는 결과로 옳은 것은? (단, 물가수준은 변하지 않는다고 가정한다.)
\begin{choices}
  \item 고정환율제도에서는 순수출이 증가한다.
  \item 고정환율제도에서는 국내 이자율이 하락한다.
  \item 변동환율제도에서는 국민소득이 변하지 않는다.
  \item 변동환율제도에서는 자국 통화가치가 절하된다.
  \item 환율제도에 관계없이 국내 통화공급량이 증가한다.
\end{choices}
\end{question}
\answer{3}
